\section{The Three Interpretations Compared}
\label{sec:comparison}

The three interpretations put different questions to the same texts and
draw chiefly on different ones: the first on Isaiah's cry and the psalms, the
second on the Gospel, the third on Isaiah's confession, the psalm's refrain
and Paul's thanksgiving.

\begin{comparisontable}{}{Rend the heavens}{The watch at the door}{The clay and the potter}
Question & Has the prophet's cry been answered, and how is it still prayed? &
Why is the hour hidden, and how is the command addressed to all? & On what
ground can a people that has strayed hope to be found blameless?\\
Decisive words & ``O that thou wouldst rend the heavens'' and, in the Latin,
``thou didst come down'' (Is 63:19b; 64:2b) & ``For ye know not when the
time is \ldots\ what I say to you, I say to all: Watch'' (Mk 13:33, 37) &
``We are clay \ldots\ thou art our father'' (Is 64:7); ``God is faithful''
(1~Cor 1:9)\\
Principal reasoning & The past tense after the plea is the answer: the Word
made flesh (Jerome); the one verb is read of the Exodus, the Incarnation and
the judgment (Aquinas); Christ, given as man, is still to be given as God
(Augustine); Christ comes twice, unseen and in glory (Cyril) & The
hour is hidden so that readiness never lapses (Jerome, Chrysostom); Christ
said ``you do not know'', not ``we'' (Jerome) and not ``I'' (Chrysostom); the
last day of the world comes to
each at his death (Augustine); the doorkeeper, the servants and every heart
each keep their watch (Bede) & God does not cause the straying (Jerome:
patience; Aquinas: permission); the anger follows the sin, and salvation is
by mercy (both); conversion itself is God's gift (Augustine); God keeps
his call (Chrysostom, Aquinas)\\
Where the Missal's own texts speak most directly & The Preface's two comings;
the solemn blessing & The Preface's \latin{vigilantes}; the Collect's works &
The Collect's petition for a will; the Prayer over the Offerings' gifts from
God's gifts\\
Where the Fathers differ & Jerome and Aquinas disagree on the literal sense
of ``thou didst come down'' & Augustine (each one's death) and Bede (the
return, divided by office) differ in emphasis on the reach of the command &
Jerome (patience) and Aquinas (permission) give two compatible accounts of why
God ``made us to err''\\
\end{comparisontable}

\subsection{Where they agree and where they press on one another}

All three take the first coming as given and the second as awaited. The
first asks about the comings themselves; the second and third ask about the
time between them, the one what it requires of those who wait, the other
what they can rely on while they wait.

They also correct one another. The first hears in Isaiah's chapter a prayer
answered; the third hears in the same chapter a confession of sin and an
appeal to fatherhood. The Latin of the reading has both. A preacher working
from the English proclaimed in the United States, which lacks the answering
clause, is closer to the third interpretation than to the first.

The second interpretation, taken alone, could sound as though the servants
secured their own readiness by effort. Against that, the third points out that the will to go
out to meet Christ is itself asked for in the Collect, and that it is Christ
who ``will confirm you unto the end''. The second in turn keeps the third
from passivity: what the master gives his servants is authority \latin{cuiusque
operis}, over every work, and Bede reads that authority as the grace of the
Holy Spirit given for good works.
